\section{Introdução}

A transição paradigmática da odontologia analógica para o ecossistema digital encontra no conceito de gêmeos digitais (\textit{Digital Twins}) sua fronteira tecnológica mais promissora. Gêmeos digitais não são meras representações tridimensionais, mas simulações dinâmicas e estocásticas das estruturas bucais e maxilofaciais do paciente, alimentadas ininterruptamente por um fluxo de dados heterogêneos. Como destacam estudos recentes de revisão escopo \cite{duggal2026exploring, katsoulakis2024digital}, a literatura científica tem testemunhado um crescimento exponencial nas publicações entre 2024 e 2026, embora o campo permaneça em um estágio de maturidade emergente, concentrando-se fundamentalmente em validações conceituais e prototipagem.

\destaque{A verdadeira disrupção tecnológica não reside apenas na virtualização da anatomia, mas na democratização do acesso a essas tecnologias por meio de plataformas \textit{open-source}.}

Historicamente, o mercado de software odontológico foi dominado por soluções proprietárias, caracterizadas por elevados custos de licenciamento e arquiteturas de caixa-preta (\textit{black-box}), as quais limitam a personalização clínica e impedem a auditoria algorítmica. Este modelo de negócio restringe severamente a pesquisa científica independente e a reprodutibilidade metodológica. Em contraponto, o movimento \textit{open-source} tem consolidado uma via alternativa robusta, fundamentada na transparência, na arquitetura composicional e na colaboração descentralizada \cite{robles2023opentwins}. Tais atributos não constituem apenas um modelo de licenciamento de software, mas uma necessidade epistemológica para a prática baseada em evidências em tempos de inteligência artificial generativa.

O ecossistema \textit{open-source} na odontologia digital transcende ferramentas isoladas. Ele abrange desde pilares consolidados de processamento de imagens médicas e modelagem geométrica (e.g., 3D Slicer, ITK, MeshLab) até infraestruturas complexas de simulação biomecânica e fluidodinâmica (e.g., FEniCS, OpenFOAM). Paralelamente a essas soluções, observamos a ascensão de arquiteturas orquestradas por ecossistemas multiagente. O OpenCode Ecosystem \cite{opencode2026ecosystem}, por exemplo, ilustra a convergência entre LLMs, repositórios abertos, motores de raciocínio formal e validação clínico-computacional contínua.

Este capítulo propõe um mapeamento crítico e sistemático destas plataformas. Em vez de uma mera catalogação de ferramentas, apresentamos uma análise de suas capacidades arquiteturais e fisiológicas para suportar a orquestração de gêmeos digitais odontológicos. Ao longo das próximas seções, delinearemos como a interoperabilidade entre agentes inteligentes e módulos matemáticos de alta precisão pode transpor a atual dependência de sistemas fechados, projetando um horizonte de pesquisa e prática clínica alinhado à medicina de precisão \cite{towards2025precision}.
